\documentclass[11pt]{article}

% ---------------------------------------------------------------
%  AI Applications in Research: The Use of Claude Code
%  A conceptual summary aimed at empirical / quantitative
%  social-science researchers.
% ---------------------------------------------------------------

\usepackage[utf8]{inputenc}
\usepackage[T1]{fontenc}
\usepackage[margin=1in]{geometry}
\usepackage{enumitem}
\usepackage{hyperref}
\usepackage{titlesec}
\usepackage{xcolor}
\usepackage{parskip}

\hypersetup{
  colorlinks=true,
  linkcolor=blue!55!black,
  urlcolor=blue!55!black,
  citecolor=blue!55!black
}

\titleformat{\section}{\large\bfseries}{\thesection}{0.6em}{}
\titleformat{\subsection}{\normalsize\bfseries}{\thesubsection}{0.6em}{}

\title{\textbf{AI Applications in Research}\\[0.3em]
       \large The Use of Claude Code in Empirical and Quantitative Social Science}
\author{Prepared as a conceptual summary}
\date{September 29, 2026}

\begin{document}
\maketitle

\begin{abstract}
\noindent
Agentic artificial intelligence is reshaping how empirical research is
conducted, moving beyond passive ``autocomplete'' assistance toward
autonomous agents that can plan, write, run, and audit code on a
researcher's behalf. This document summarizes the main concepts behind
AI applications in research, with particular emphasis on
\emph{Claude Code}, a command-line coding agent, and how it is being
adopted for quantitative and empirical social-science workflows. The
discussion is framed around themes popularized in ongoing practitioner
writing---most visibly Scott Cunningham's ``Scott's Mixtape'' Substack
series on Claude Code for social scientists.
\end{abstract}

\section{Introduction}

Research computing has historically relied on the researcher writing every
line of analysis code by hand. The emergence of large language models
(LLMs) first introduced \emph{in-editor autocomplete} and chat-based
assistance. The current frontier is \emph{agentic AI}: systems that
operate an environment---reading files, executing commands, inspecting
output, and iterating---under a researcher's supervision. Claude Code is
one such agent that runs in the terminal (and in web, desktop, and IDE
contexts) and can carry a multi-step empirical task from specification to
result.

For empirical social scientists---economists, political scientists,
sociologists, epidemiologists---this shift matters because the bottleneck
in applied work is rarely the idea; it is the labor of data cleaning,
estimation, robustness checking, and replication. Agentic tools promise to
compress that labor while, if used carefully, \emph{raising} the standard
of transparency and reproducibility.

\section{Core Concepts}

\subsection{From Autocomplete to Agents}
An \textbf{AI agent} differs from a chatbot in that it can take actions in
a real environment and observe their consequences. Rather than returning a
block of suggested code, the agent runs the code, reads the error, and
fixes it---looping until a goal is met. This ``plan--act--observe'' loop is
the defining feature of tools such as Claude Code.

\subsection{The Agent Harness}
A recurring theme is the idea of an \textbf{agent harness}: the scaffolding
of instructions, project context, permissions, and repository conventions
that steers a general-purpose agent toward reliable, domain-appropriate
behavior. In a research setting this includes:
\begin{itemize}[leftmargin=1.4em]
  \item A project-level context file (e.g.\ a \texttt{CLAUDE.md}) describing
        the dataset, the estimands, coding conventions, and ``house rules.''
  \item Explicit statements of the empirical design (identification
        strategy, sample restrictions, fixed effects, clustering).
  \item Guardrails on what the agent may and may not do (e.g.\ never alter
        raw data, always log every transformation).
\end{itemize}
A good harness turns a generic coding agent into a disciplined research
assistant.

\subsection{Reproducibility and Auditing}
Because the agent produces a full, inspectable transcript of commands and
edits, agentic workflows can \emph{improve} reproducibility. A prominent
use case is the \textbf{code audit}: asking the agent to reproduce a
published or draft result and to flag discrepancies. Because agents are
increasingly language-agnostic, one audit can span \textbf{Python, R, and
Stata}, checking that the same empirical estimator (for example, a
difference-in-differences design) yields consistent numbers across
implementations.

\section{Applications in Empirical Research}

\begin{enumerate}[leftmargin=1.5em]
  \item \textbf{Data cleaning and wrangling.} Automating the tedious,
        error-prone steps of merging, reshaping, and validating data, with
        the agent documenting each transformation.
  \item \textbf{Estimation and modeling.} Implementing econometric
        estimators---panel fixed effects, instrumental variables, modern
        difference-in-differences (e.g.\ Callaway--Sant'Anna style),
        synthetic control---and generating the accompanying tables and
        figures.
  \item \textbf{Cross-language replication audits.} Re-implementing an
        analysis in a second language to detect coding errors or
        package-specific default differences.
  \item \textbf{Robustness and specification search.} Systematically
        varying controls, samples, and functional forms to map how
        sensitive a result is---made explicit and logged rather than
        hidden.
  \item \textbf{Literate documentation.} Producing readable write-ups,
        replication packages, and README files so that a project is
        legible to co-authors and referees.
  \item \textbf{Live and exploratory research.} Using the agent as a
        pair-programmer during active projects, documenting ``what works
        and what doesn't'' in real time.
\end{enumerate}

\section{Benefits and Risks}

\subsection{Benefits}
\begin{itemize}[leftmargin=1.4em]
  \item \textbf{Speed.} Compresses routine implementation work from hours to
        minutes.
  \item \textbf{Breadth.} Lowers the cost of working across unfamiliar
        languages and packages.
  \item \textbf{Transparency.} Every step is captured in a transcript that
        can be reviewed and re-run.
  \item \textbf{Learning.} Acts as an always-available tutor for methods and
        tooling.
\end{itemize}

\subsection{Risks and Mitigations}
\begin{itemize}[leftmargin=1.4em]
  \item \textbf{Silent errors / hallucination.} An agent can produce
        plausible-but-wrong code. \emph{Mitigation:} the researcher remains
        the accountable reviewer; results are cross-checked and audited.
  \item \textbf{Over-delegation.} Losing understanding of one's own
        pipeline. \emph{Mitigation:} treat the agent as an assistant, not an
        author of record; keep human judgment on identification and
        interpretation.
  \item \textbf{Reproducibility of the agent itself.} Model behavior can
        drift. \emph{Mitigation:} commit the generated code and data
        pipeline, not just the prompts.
  \item \textbf{Data governance.} Sensitive or restricted data require
        care about where computation runs and what is shared.
\end{itemize}

\section{Practical Workflow Pattern}

A representative agentic research loop:
\begin{enumerate}[leftmargin=1.5em]
  \item \textbf{Specify.} State the empirical goal, design, and constraints
        in a project context file.
  \item \textbf{Delegate.} Ask the agent to implement a discrete,
        well-scoped step.
  \item \textbf{Inspect.} Read the diff, the logs, and the output; question
        anything surprising.
  \item \textbf{Audit.} Have the agent (or a second pass) reproduce or
        cross-check the result, ideally in another language.
  \item \textbf{Commit.} Version-control the code, outputs, and
        documentation so the work is durable and shareable.
\end{enumerate}

\section{Connection to Current Practitioner Writing}

Much of the practical knowledge about using Claude Code in research is being
developed in public. Scott Cunningham's ``Scott's Mixtape'' Substack runs an
ongoing series documenting how agentic AI is changing empirical
social-science work---covering agent harnesses for economics, cross-language
diff-in-differences audits, and hands-on tutorials on using Claude Code for
quantitative research. Such practitioner documentation is currently one of
the best guides to \emph{what works and what does not} in day-to-day use.

\section{Conclusion}

AI applications in research have moved from suggestion to action. Claude
Code and similar agents can meaningfully accelerate empirical work while,
used responsibly, strengthening transparency and reproducibility. The
central principle is unchanged: the researcher owns the identification
strategy, the interpretation, and the accountability---the agent is a
powerful, auditable assistant, not a replacement for scientific judgment.

\vspace{1em}
\hrule
\vspace{0.6em}
\footnotesize
\noindent\textbf{Key references / further reading:}
\begin{itemize}[leftmargin=1.4em, itemsep=0.1em]
  \item Scott's Mixtape Substack --- Claude Code series:
        \url{https://causalinf.substack.com/s/claude-code}
  \item ``Claude Code 51: Harnessing AI Agents for Economic Research'':
        \url{https://causalinf.substack.com/p/claude-code-51-harnessing-ai-agents}
  \item ``Claude Code Part 12: How I Use Claude Code for Empirical
        Research'':
        \url{https://causalinf.substack.com/p/claude-code-part-12-how-i-use-claude}
  \item Claude Code documentation:
        \url{https://code.claude.com/docs}
\end{itemize}

\end{document}
